\documentclass{article}
\usepackage{amsmath}
\usepackage{amssymb}
\usepackage{booktabs}
\usepackage[a4paper,margin=1in]{geometry}

\title{Model Training Recipes (T18--T20)\\ \large Reference and plain-language explanation}
\date{July 18, 2026}
\author{}

\begin{document}
\maketitle

\begin{abstract}
This note records the exact training recipe for the three ensemble members --- ProkBERT (T18), Nucleotide Transformer + LoRA (T19), and the from-scratch CNN (T20) --- and explains, in plain language, what every hyperparameter and component means and why it is set the way it is. The values below are the standalone-run defaults that produced the current verified models; they have not yet been changed by the T21 hyperparameter sweep.
\end{abstract}

\section{The three models at a glance}
All three are binary classifiers that read an $81$\,bp DNA window (produced by T16) and output a promoter / not-promoter score. Two are \emph{fine-tuned} from a pretrained model (ProkBERT, NT); one is \emph{trained from scratch} (CNN).

\begin{table}[h]
\centering
\caption{Training recipe for the three ensemble members (standalone-run defaults).}
\begin{tabular}{lccc}
\toprule
Setting & ProkBERT (T18) & NT-LoRA (T19) & CNN (T20) \\
\midrule
Base / type        & \texttt{prokbert-mini} (full FT) & NT v2 250M (LoRA) & from scratch \\
Learning rate      & $2\times10^{-5}$ & $1\times10^{-4}$ & $1\times10^{-3}$ \\
Batch size         & 16 & 8 & 64 \\
Max epochs         & 10 & 10 & 30 \\
Warmup             & ratio 0.1 & 200 steps & none \\
Weight decay       & 0.01 & 0.01 & $1\times10^{-5}$ \\
Dropout            & 0.1 & 0.1 (LoRA) & 0.3 \\
Optimizer          & AdamW & AdamW & AdamW \\
LR schedule        & linear & linear & cosine warm restarts \\
Early-stop patience& 3 & 3 & 8 \\
Precision          & bf16 & bf16 & fp32 \\
\midrule
Val AUPRC (result) & 0.876 & 0.853 & 0.811 \\
\bottomrule
\end{tabular}
\end{table}

\section{Common ingredients (shared by all three)}
These terms appear for every model; they are explained once here.

\begin{description}
\item[Learning rate (LR).] The step size for each weight update --- how far the model moves in the direction the gradient points. Too high and training overshoots or diverges; too low and it learns slowly or gets stuck. Note the pattern in the table: the two \emph{fine-tuned} models use tiny LRs ($2\times10^{-5}$, $1\times10^{-4}$) so they nudge their already-good pretrained weights without destroying them, while the \emph{from-scratch} CNN uses a much larger LR ($1\times10^{-3}$) because its weights start random and need to move a long way.

\item[Batch size.] How many $81$\,bp windows are processed together before one weight update. Bigger batches give a smoother, more reliable gradient estimate but use more GPU memory; smaller batches are noisier (which can even help generalization). The 250M-parameter NT model uses the smallest batch (8) because it is the biggest and least fits in memory; the tiny CNN uses 64.

\item[Epoch / max epochs.] One epoch is one full pass over the whole training set. ``Max 10'' (or 30) is only a ceiling --- early stopping (below) almost always halts training earlier.

\item[Warmup.] At the very start of training, the learning rate is ramped up gradually from near-zero to its full value, rather than starting at full strength. \texttt{ratio 0.1} means this ramp lasts the first $10\%$ of training steps; \texttt{200 steps} means the first 200 updates. Warmup prevents large, destabilizing updates before the optimizer's internal statistics have settled --- important for transformers.

\item[Weight decay.] A regularization term that gently pulls every weight toward zero on each step (an $L_2$ penalty). It discourages the model from relying on a few very large weights, which reduces overfitting. $0.01$ is a standard value for transformers; the CNN uses a much smaller $1\times10^{-5}$.

\item[Dropout.] During training only, a fraction of the internal activations is randomly set to zero on each step ($0.1$ = $10\%$, $0.3$ = $30\%$). This forces the network not to depend on any single feature and is a strong anti-overfitting tool. It is switched off at inference. The from-scratch CNN uses a higher rate (0.3) because it overfits more easily than a pretrained transformer.

\item[Optimizer: AdamW.] The algorithm that converts gradients into actual weight updates. ``Adam'' adapts the step size separately for each parameter using running averages of the gradient's mean and variance (so it needs little manual tuning); the ``W'' means weight decay is applied correctly and separately from the gradient. All three models use it.

\item[LR schedule.] How the learning rate changes \emph{over} training (after any warmup). \textbf{Linear} (the two transformers) decays the LR in a straight line down to zero by the final step. \textbf{Cosine warm restarts} (the CNN) is explained in the CNN section.

\item[Early-stopping patience.] After each epoch the model is scored on the validation set; if the chosen metric has not improved for \texttt{patience} consecutive epochs, training stops and the \emph{best} checkpoint seen so far is kept. Patience 3 suits the transformers (they converge in a few epochs); the CNN gets 8 because it trains longer. This prevents both overfitting and wasted compute.

\item[Model-selection metric: AUPRC.] The validation metric used both to pick the best epoch and to trigger early stopping is the area under the precision--recall curve (\texttt{eval\_auprc}). AUPRC is chosen (over accuracy or AUROC) because the eventual genome task is highly imbalanced, and AUPRC is the metric that reflects performance on the rare positive class.

\item[Mixed precision (bf16 / fp32).] ``bf16'' means the math is done in 16-bit ``brain floating point'' instead of 32-bit, which is roughly twice as fast and uses half the memory on the H800 GPU; bfloat16 keeps the wide numeric range of 32-bit (unlike fp16), so it is numerically stable on Hopper GPUs. The CNN is small enough that it just stays in full 32-bit precision.

\item[The weighted loss.] Every training window carries a number, \texttt{weight}. Instead of averaging the per-window loss equally, all three models minimise a \emph{weighted} average:
\begin{equation}
\mathcal{L} = \frac{\sum_i w_i \cdot \mathrm{CE}\!\left(\hat{y}_i, y_i\right)}{\sum_i w_i},
\end{equation}
where $\mathrm{CE}$ is the cross-entropy loss for window $i$ and $w_i$ is its weight. The weight has two jobs multiplied together: (i) \emph{label confidence} --- ``silver'' positives (the inferred T13 housekeeping upstreams) get $\times 0.3$ so they influence training less than verified promoters; and (ii) \emph{class balance} --- the negatives are rescaled so the effective training set is $\sim$50/50 positive/negative. Validation and test are left \emph{unweighted}, at the natural $\sim$29\% positive rate, so the reported metrics are honest.

\item[Infrastructure (preempt + checkpoint).] All jobs run on the cheap preemptible SLURM queue; a signal trap (SIGUSR1) saves the model just before the job can be evicted, so it resumes rather than restarts. \texttt{seed=42} fixes the randomness for reproducibility.
\end{description}

\section{T18 --- ProkBERT (full fine-tune)}
ProkBERT-mini is a small BERT-style transformer pretrained on prokaryotic DNA. Here it is \emph{fully} fine-tuned (all weights update) on the T16 windows, with a fresh 2-class classification head. The recipe is the textbook small-BERT fine-tune: LR $2\times10^{-5}$, batch 16, up to 10 epochs, linear warmup over the first $10\%$ of steps, weight decay 0.01, dropout 0.1, bf16, early stopping (patience 3) on validation AUPRC, keeping the best checkpoint. It uses the shared weighted loss above. Result: validation AUPRC $0.876$ (test $0.872$) --- the strongest of the three.

\section{T19 --- Nucleotide Transformer + LoRA}
NT v2 250M is a much larger transformer (250 million parameters) pretrained across many species. Fully fine-tuning it would be expensive and prone to overfitting on our modest data, so it is adapted with \textbf{LoRA (Low-Rank Adaptation)}: the 250M pretrained weights are \emph{frozen}, and small trainable matrices are inserted into a few layers. Only about $0.5\%$ of the parameters actually train.

\begin{description}
\item[\texttt{r} = 8 (rank).] The size of the low-rank update inserted into each targeted layer --- effectively a bottleneck dimension. Higher \texttt{r} gives the adapter more capacity but more parameters to train; 8 is small and cheap.
\item[\texttt{lora\_alpha} = 16 (scaling).] A scale factor applied to the LoRA update. The effective scale is $\alpha / r = 16/8 = 2$; it controls how strongly the adapter's output is added to the frozen layer.
\item[\texttt{lora\_dropout} = 0.1.] Dropout applied inside the LoRA layers, for the same anti-overfitting reason as ordinary dropout.
\item[\texttt{target\_modules} = [query, value].] \emph{Which} weight matrices receive an adapter --- here the ``query'' and ``value'' projections of the attention mechanism, the standard and most effective choice for this kind of model.
\item[\texttt{modules\_to\_save} = [classifier].] The brand-new classification head is \emph{not} frozen --- it is trained in full, because it starts from random initialization and must be learned from scratch.
\item[\texttt{task\_type} = SEQ\_CLS.] Tells the PEFT library this is a sequence-classification task, so it wires the adapter and head correctly.
\end{description}

\noindent Its optimisation settings: LR $1\times10^{-4}$ (a bit higher than ProkBERT, typical for LoRA, since only small adapters move), batch 8, up to 10 epochs, 200 warmup steps, weight decay 0.01, bf16, early stopping (patience 3) on AUPRC. \textbf{Gradient checkpointing} is available (it recomputes activations during back-propagation instead of storing them, trading extra compute for much less memory so the 250M model fits on one GPU); the verified run used \texttt{-{}-no-grad-checkpoint} because it fit without it. It runs on the separate \texttt{nt-tf4x} software environment (transformers 4.57), because NT's model code is incompatible with the newer transformers 5.x that the other two use. Result: validation AUPRC $0.853$ (test $0.841$), with only $0.507\%$ of parameters trainable.

\section{T20 --- Custom CNN (from scratch)}
The CNN is the ``control'': a small convolutional network trained from random initialization, no pretraining. It reads the $81\times4$ one-hot encoding of the window (4 channels, one per DNA base).

\begin{description}
\item[Three Conv1d blocks (64 $\to$ 128 $\to$ 256 filters).] Each 1D convolution slides a set of small learnable filters along the sequence, detecting local patterns (the building blocks of a promoter motif). Deeper blocks use more filters, so the network can recognise more distinct patterns as it builds up.
\item[Kernel sizes 7 / 5 / 3.] The width (in bp) of the convolution window in each block --- a 7\,bp view first, then narrower 5 and 3\,bp views.
\item[BatchNorm.] Normalises the activations within each batch, which stabilises and speeds up training. On by default (\texttt{-{}-no-batch-norm} disables it).
\item[MaxPool.] After the first two blocks, downsampling that keeps only the strongest activation in each small window --- it shrinks the sequence length and makes the detector tolerant to a motif shifting position.
\item[Dropout 0.3.] Regularisation (higher than the transformers, since a from-scratch model overfits more readily).
\item[AdaptiveMaxPool1d(1) $\to$ FC 128 $\to$ 2.] A global max-pool takes the single strongest response of each filter across the whole sequence, giving a fixed-length vector; a fully-connected layer maps it to 128 hidden units and then to 2 output logits (promoter vs.\ not).
\end{description}

\noindent Optimisation: AdamW, LR $1\times10^{-3}$, weight decay $1\times10^{-5}$, batch 64, up to 30 epochs, early stopping (patience 8) on AUPRC, and \textbf{gradient clipping} at \texttt{max\_norm} $=1.0$ (caps the size of any single update to prevent occasional ``exploding'' gradients). Its schedule is \textbf{CosineAnnealingWarmRestarts} (\texttt{T\_0}=5, \texttt{T\_mult}=2): the LR follows a cosine curve down over 5 epochs, then jumps back up (``restarts'') and repeats over ever-longer cycles ($5, 10, 20, \dots$), which helps a small network escape shallow local minima. (A known caveat: these restarts interact awkwardly with patience-based early stopping and are worth revisiting in the T21 sweep.) Result: validation AUPRC $0.811$ (test $0.818$).

\section{T21 --- hyperparameter sweep (not yet run)}
The values above are sensible defaults, not tuned optima. T21 is an Optuna search that will vary each model's key knobs --- learning rate, batch size, warmup, weight decay, dropout, epochs (and for NT the LoRA rank/$\alpha$) --- and keep whichever combination maximises validation AUPRC. Until it runs and the winners are retrained, the recipe above is what the current models use.

\end{document}
